求比值、极值和相对位置的分布，可从联合密度作变量变换，写清取值域与雅可比。条件缩放有时能消去已知位置，直接证明独立性。均匀切点产生的间距在单纯形上具有常密度，但各间距并不独立。

几何概率必须指定位置与方向的抽样方式。距离阈值可转为邻域内的二项计数，交点与配对可用指示变量求均值；共点或边界情形须单独辨认。分布对称也不保证期望存在。

随机着色可用非负整数坏事件数的期望小于一证明存在性。保密问题则要比较给定消息的密文分布，不能仅凭密文边际均匀判断联合保密。
\par\smallskip
